\documentclass[12pt]{article}
\usepackage[utf8]{inputenc}
\usepackage[T1]{fontenc}
\usepackage{hyperref}
\title{Wali Language Resources for K-12 Curriculum Development}
\author{Compiled by Claude Sonnet 5.5 (claude-sonnet-5-5) for Andr\'as Kornai}
\date{}
\begin{document}
\maketitle

\section{Introduction}
Wali (ISO 639-3 wll) is a Hill Nubian language of the Nuba Mountains, Sudan. Speaker estimates vary by source. Wikipedia and Ethnologue give about 9,000, and Joshua Project cites 2,700. UNESCO classes it as critically endangered. Very little published material exists.

\section{Primary Educational Resources}
No primers, textbooks or school programmes were found in search.

\section{Linguistic and Reference Materials}
\begin{itemize}
\item ``An initial sketch on Wali phonology (a Sudanese endangered language)'', \emph{Adab} (University of Khartoum journal). It reports 19 consonants and 8 vowels, based on two informants. \url{https://onlinejournals.uofk.edu/JFA/article/view/2083}
\item Wikipedia, Wali language (Sudan): \url{https://en.wikipedia.org/wiki/Wali_language_(Sudan)}
\item Glottolog entry: \url{https://glottolog.org/resource/languoid/id/wali1262}
\item ASJP word list: \url{https://asjp.clld.org/languages/WALI_SUDAN}
\item Verbix endangered-languages page, which gives the autonym and related terms: \url{https://docs.verbix.com/EndangeredLanguages/Wali}
\item Krell (2012), SIL sociolinguistic survey of the Ama, Karko and Wali groups. It is mentioned in search results. I saw no URL.
\end{itemize}

\section{Religious and Community Resources}
Joshua Project describes the Wali as a small Nuba people and lists no written materials. It gives Bible translation status as ``translation started''. \url{https://www.joshuaproject.net/index.php/people_groups/15836}

\section{Digital and Audio Resources}
None found.

\section{Curriculum Development Context}
No dictionary, grammar or learner materials were found. A curriculum would need new documentation work with the community, building on the phonology sketch and the SIL survey. Resources are very limited.
\end{document}